%=============================================================================
% Motivation und Zielsetzung (eigenstaendiges Dokument)
% Seminararbeit IT-Sicherheit, M.Sc. IVS, Hochschule Bremerhaven
% Autor: Jordan Guimfack Jeuna (38184)
% Kompilierung: pdflatex + bibtex
%=============================================================================
\documentclass[letterpaper,twocolumn,10pt]{article}
\usepackage{usenix-2020-09}

% Sprache und Zeichenkodierung
\usepackage[utf8]{inputenc}
\usepackage[T1]{fontenc}
\usepackage[ngerman]{babel}

% Typografie und Mathematik
\usepackage{microtype}
\usepackage{amsmath}
\usepackage{amssymb}

% Hyperlinks (Bookmarks deaktiviert)
\hypersetup{bookmarks=false,colorlinks=true,linkcolor=black,citecolor=black,urlcolor=black}

% Abkuerzungen
\newcommand{\shap}{SHAP}
\newcommand{\pca}{PCA}
\newcommand{\ba}{Balanced Accuracy}
\newcommand{\gb}{Gradient Boosting}
\newcommand{\iot}{IoT}
\newcommand{\ids}{IDS}

%=============================================================================
\begin{document}
\date{}

\title{\Large \bf
  Erklaerbare Erkennung von Cyberangriffen in IoT-Netzwerken:\\
  Entwicklung und Evaluation eines Gradient-Boosting-Klassifikators\\
  mit SHAP-Analyse auf dem CICIoT2023-Datensatz\\[6pt]
  \large Motivation und Zielsetzung
}

\author{
  {\rm Jordan Guimfack Jeuna}\ (Matrikelnummer 38184)\\
  Hochschule Bremerhaven\\
  Studiengang M.Sc. Informatik (Vertrauenswuerdige Systeme)\\
  Modul IT-Sicherheit, Prof.\ Dr.\ Lars Fischer
}

\maketitle

%=============================================================================
\section{Motivation}\label{motivation}

Der CICIoT2023-Datensatz~\cite{neto2023ciciot2023} erfasst 33 reale Angriffe gegen
ein Testnetz aus 105 \iot{}-Geraeten und macht damit messbar, wie breit die
Angriffsflaeche vernetzter Geraete inzwischen ist.
Signaturbasierte \ids{} stossen in diesem Umfeld an eine klare Grenze: Sie erkennen
nur Angriffe, deren Muster bereits hinterlegt ist, und versagen bei zuvor unbekannten
Varianten.
Machine Learning loest dieses Problem datengetrieben, indem es Angriffsverhalten aus
Netzwerkmerkmalen lernt statt aus festen Signaturen.
Baumbasierte Ensemble-Verfahren bilden dabei den belastbaren Stand der
Forschung~\cite{kikissagbe2024review}, und \gb{} erreicht auf dem CICIoT2023-Datensatz
die hoechste \ba{} unter den in Raturi et al.\ verglichenen
Modellen~\cite{raturi2026exploratory}.

Diese Leistung allein genuegt fuer den praktischen Einsatz jedoch nicht.
Die einschlaegigen Arbeiten auf CICIoT2023 werten ihre Modelle als Black Box aus und
legen nicht offen, welche Netzwerkmerkmale eine Klassifikation
tragen~\cite{raturi2026exploratory,almahaqeri2026gradient}.
Fuer ein Sicherheitssystem ist genau das entscheidend: Ein Analyst, der einen Alarm
prueft, muss erkennen koennen, warum ein Paket als Angriff eingestuft wurde, bevor er
darauf reagiert.

%=============================================================================
\section{Forschungsluecke}\label{luecke}

Aus dem aktuellen Stand ergeben sich drei konkrete Schwachstellen, an denen diese
Arbeit ansetzt.

\textbf{Verzerrte Bewertung durch Klassenungleichgewicht.}
DoS- und DDoS-Verkehr dominiert den CICIoT2023-Datensatz zahlenmaessig.
Ein Modell kann seltene Angriffsklassen vollstaendig verfehlen und trotzdem eine hohe
Accuracy ausweisen, weil die haeufigen Klassen das Ergebnis bestimmen.
Hosseini et al.~\cite{hosseini2025imbalance} belegen diese Verzerrung fuer
\ids{}-Datensaetze und machen die \ba{} als Bewertungsmassstab erforderlich.

\textbf{Fehlende Interpretierbarkeit.}
Ohne Einblick in die Entscheidungsgrundlage bleibt unklar, ob ein Modell ein echtes
Angriffsmerkmal nutzt oder einem Artefakt des Datensatzes folgt.
Keine der genannten Arbeiten auf CICIoT2023 liefert diese Begruendung pro
Angriffsklasse.

\textbf{Unbewerteter Einfluss der Dimensionsreduktion.}
Raturi et al.~\cite{raturi2026exploratory} reduzieren die 46 Merkmale per \pca{} auf
16 Komponenten, ohne zu pruefen, wie sich dieser Schritt auf die Erkennungsleistung
und auf die Nachvollziehbarkeit der Merkmale auswirkt.
Da \pca{}-Komponenten Linearkombinationen der Originalmerkmale sind, steht zu
erwarten, dass die fachliche Deutbarkeit darunter leidet.

%=============================================================================
\section{Zielsetzung}\label{ziel}

Die Arbeit entwickelt einen \gb{}-Klassifikator auf dem CICIoT2023-Datensatz, der
zwei Anforderungen zugleich erfuellt: eine hohe Erkennungsleistung, gemessen an der
\ba{} als primaerer Metrik, und eine pro Angriffsklasse nachvollziehbare
Entscheidungsgrundlage durch \shap{}.
Eine Ablation Study stellt eine Pipeline mit 16 \pca{}-Komponenten einer
Kontrollpipeline mit allen Originalmerkmalen gegenueber und beziffert so den
Preis der Dimensionsreduktion fuer Leistung und Erklaerbarkeit.

Der eigenstaendige Beitrag liegt in der Verbindung dreier bislang getrennter Elemente
auf CICIoT2023: \shap{} als Erklaerungskomponente, die \ba{} als primaerer Massstab
und eine messbare Pruefung der Erklaerungen.
Die Qualitaet der \shap{}-Ausgaben wird dabei nicht unterstellt, sondern an den Kriterien
Fidelitaet, Konsistenz und Jaccard-Aehnlichkeit nach Hermosilla et
al.~\cite{hermosilla2025xai_forensic} ausgerichtet und qualitativ gegen die bekannten
Netzwerksignaturen der jeweiligen Angriffstypen geprueft.
Die Wahl von \shap{} stuetzt sich auf das PRISMA-Review von Ogunseyi et
al.~\cite{ogunseyi2026xai_review}, das \shap{} in der Mehrheit von 129 untersuchten
XAI-\ids{}-Studien einsetzt, ohne dass die Erkennungsleistung darunter leidet.

% Literaturverzeichnis
{\footnotesize \bibliographystyle{plain}
\bibliography{references}}

\end{document}
